\section{Algorithm Correctness}
\label{app:algorithm-proofs}

For Sections~\ref{app:proof-sw}--\ref{app:proof-limit}, fix a strict total order $\prec$ over row identities under instruction $\phi$ with deterministic tie-breaking. A \emph{consistent oracle} returns the restriction of $\prec$ to any batch $B$. Section~\ref{app:proof-inconsistent} requires only that each validated call return a permutation of $B$; distinct input records always retain distinct identities.

\subsection{Sliding-Window Merge}
\label{app:proof-sw}

\begin{lemma}[Safe window prefix]
\label{lem:sw-safe}
For $\prec$-sorted runs and a consistent oracle, each round of Algorithm~\ref{alg:sliding_window_kway_group_merge} emits a nonempty prefix of the $\prec$-sorted union of unemitted rows whenever rows remain, ranking at most $|A|h \le m$ rows (adaptive) or $gh \le kh \le m$ rows (fixed-quota $\mathrm{SW}(k,h)$ on $g \le k$ runs).
\end{lemma}
\begin{proof}
Maintain the invariant that after refill, each active run $u$'s buffered slice $W_u$ is a contiguous prefix of its remaining rows. This holds at $p_u = 0$ and is preserved when removing a prefix per run. If $|A|=1$, the lone run's sorted suffix is appended without a call. Otherwise, for each $u \in A$ with unread items ($p_u + b_u < |R_u|$), let $z_u$ be its last buffered item ($b_u = h \ge 1$) at index $j_u$ in $O$. Emission stops at $t = \min_u j_u$, or drains $O$ if no unread suffix exists.

Every emitted item $x$ precedes or equals $z_u$ in $O$ (hence in $\prec$) for all such $u$. Since unread items of $u$ follow $z_u$ in $\prec$, $x$ precedes all unread suffixes and all survivors $O[t{+}1:]$. Thus the emitted slice is the exact next prefix of the sorted union, and its restriction to each run is a prefix of $W_u$, restoring the invariant. Because $t \ge 1$ (or $|O| \ge 1$), every round makes progress within capacity $m$.
\end{proof}

\subsection{Segmented Loser-Tree Merge}
\label{app:proof-lt}

At internal node $v$, child streams deliver $h$-row blocks (possibly short at exhaustion) into slots $L$ and $R$. Define a slot's \emph{frontier} as its maximum real row under $\prec$ if full ($h$ rows), and $+\infty$ otherwise; retained losers may mix rows from both children.

\begin{lemma}[Node replay and frontier invariant]
\label{lem:lt-node}
Suppose $v$'s slots partition all fetched unforwarded rows, unfetched child rows follow their slot's frontier in $\prec$, and frontier $+\infty$ implies an exhausted child. Algorithm~\ref{alg:replay_segment_node} emits the next block in $\prec$ and restores the invariant once $\mathit{Pointer}(v)$ supplies its next block.
\end{lemma}
\begin{proof}
Initially each slot holds its own child's first block, satisfying the frontier invariant. Let $b_L, b_R$ be the slot frontiers when $L \uplus R$ sorts into $Z$. Any unexhausted child has $h$ slot entries $\preceq$ its frontier, so $Z[1..h]$ is $\preceq \min(b_L, b_R)$ and precedes all unfetched rows; if $|Z| < h$, $b_L = b_R = +\infty$ and emitting $Z$ drains both children.

With virtual $+\infty$ tail padding, the first slot reaching $h$ entries in $Z$ (lines~2--\ref{line:replay-split}) attains $\min(b_L, b_R)$. Meanwhile $\mathit{Loser}(v) = Z[h{+}1..2h]$ retains all unforwarded real rows: if both slots were full, $|\mathit{Loser}(v)| = h$ with maximum $\max(b_L, b_R)$; if either was short, $|\mathit{Loser}(v)| < h$ and the unpointed slot has frontier $+\infty$. Assigning $\mathit{Loser}(v)$ to the unpointed slot preserves frontier $\max(b_L, b_R)$, and refilling from $\mathit{Pointer}(v)$ restores the pointed slot's frontier. For $h{=}2$ with $[1,100]$ and $[3,4]$, replay emits $[1,3]$, retains $[4,100]$ (frontier $100$), and refills the second slot.
\end{proof}

\begin{lemma}[Halving-tree construction and path refresh]
\label{lem:lt-tree}
Algorithms~\ref{alg:loser_tree_kway_group_merge} and~\ref{alg:refresh_loser_tree_path} maintain Lemma~\ref{lem:lt-node} across the halving tree $T$, returning the $\prec$-sorted union within capacity $m$ when $(1+\lceil\log_2 g\rceil)h \le m$.
\end{lemma}
\begin{proof}
View each internal node as a block producer that initializes from its two child blocks and thereafter pulls solely along $\mathit{Pointer}(v)$. Descending from $\textsc{Root}(T)$ traces one path $\mathcal{P}$ ($|\mathcal{P}| \le d = \lceil\log_2 g\rceil$), pulls one leaf block, and replays $\mathcal{P}$ bottom-up while off-path nodes retain valid state. Under a consistent oracle, restricting a packed-init or path-refresh rank map $\rho$ to $L \uplus R$ matches ranking $L \uplus R$ directly, and run-order projection is an identity map; induction on height yields the sorted union (singleton groups pass through without a tree).

Halving $g \ge 2$ leaves gives depth $d = \lceil\log_2 g\rceil$. Since $(d+1)h \le m$ implies $\lfloor m/h \rfloor \ge 2$, each packed initialization cut contracts at least one sibling pair within width $m$, and each refresh ranks at most $(d+1)h \le m$ real rows. At $h{=}1$, $\mathit{Pointer}(v)$ tracks the single winner's source child, recovering the classical loser tree.
\end{proof}

\subsection{Materialization and Demand-Driven Prefix Execution}
\label{app:proof-limit}

\begin{theorem}[Full sortedness and \LIMIT~$K$ equivalence]
\label{thm:full-limit}
Under a consistent oracle, Algorithm~\ref{alg:semantic_external_merge} returns the $\prec$-sorted permutation $\Pi$, and the producer tree (Section~\ref{sec:topk-exec}) returns its first $\min(K,n)$ rows $\Pi_K$.
\end{theorem}
\begin{proof}
Phase~1 sorts disjoint chunks under $\prec$. By Lemmas~\ref{lem:sw-safe} and~\ref{lem:lt-tree} plus singleton pass-through, pass $t$ produces $r_t = \lceil r / k^t \rceil$ disjoint sorted runs, leaving $\Pi$ after $\lceil \log_k r \rceil$ passes ($r \ge 1$). For \LIMIT~$K$, fix initial runs and $\theta$: by induction on tree height, each child's $\textsc{Next}(q)$ delivers the exact next $q$ items of its materialized stream. Every parent refill thus sees identical inputs and state transitions as in Full, buffering unconsumed tails locally and stopping once the root returns $\min(K,n)$ items—even without global consistency whenever corresponding producer calls return matching permutations.
\end{proof}

\subsection{Guarantees under Inconsistent Rank Calls}
\label{app:proof-inconsistent}

\begin{proposition}[Row-identity and run-prefix guarantees]
\label{prop:inconsistent}
If every call returns a batch permutation, both SW and LT($h$) conserve row identities across Full and \LIMIT~$K$. With run-order projection, every LT($h$) output prefix restricts to an exact prefix of each Phase-1 run in original order; unprojected SW need not preserve run prefixes.
\end{proposition}
\begin{proof}
SW removes only emitted identities ($\ge 1$ per round), and LT replay partitions $L \uplus R$ disjointly. If one LT slot is short and one is full, real entries precede $+\infty$ padding, pointing $\mathit{Pointer}(v)$ to the full child so an empty root winner guarantees complete drainage. For any run $R$ in LT, all $R$-rows enter node $v$ via one child contiguously in run order; run-order projection admits only each run's earliest remaining item into $\rho$, so $Z[1..h]$ forwards only the next contiguous prefix of $R$ across all passes and root truncations. In SW with $R_1=(a_1,a_2,a_3,a_4)$, $R_2=(b_1,b_2)$, $h{=}2$, ranking $B$ as $(a_2,b_1,a_1,b_2)$ emits through $a_1$, so \LIMIT~$2$ outputs $(a_2,b_1)$ without $a_1$.
\end{proof}

\section{Additional Cost-Model Analysis}
\label{app:cost-results}

\subsection{Exact Expectation for Single-Item LT Prefixes}
\label{app:cost-exact}

For positive run lengths $\boldsymbol{\ell}$, $h{=}1$, and $1 \le K \le \min_i \ell_i$, packed initialization installs all run heads and emits $\pi_1$ at fixed cost $(A_0, S_0)$. Without prefetching, step $j \in \{2,\dots,K\}$ replaces only $\pi_{j-1}$: since fewer than $K \le \min_i \ell_i$ rows have left each run, source run $i$ has a replacement and full path opponents, incurring deterministic path work $(a_i, s_i)$. Under uniform rank partitioning, rank $j{-}1$ comes from run $i$ with marginal probability $\ell_i / n$; linearity of expectation yields Eq.~\eqref{eq:cost-limit-single} without assuming independent winners. Once $K > \min_i \ell_i$ or $h > 1$, finite replay captures run exhaustion and segment boundaries.

\subsection{Rank-Skew Sensitivity of Sliding-Window Prefixes}
\label{app:cost-skew}

If the top $h$ window ranks fall in one unexhausted run, SW stops at $h$ outputs; balanced interleaving extends the safe prefix. Strong rank clustering can thus cause the uniform model to overestimate safe prefixes on short limits. Because later windows inherit surviving ranks, our simulator retains assigned integer ranks across rounds and passes rather than resampling; more Monte Carlo trials reduce sampling variance under uniform interleaving but do not remove prior mismatch on skewed workloads.

\subsection{First-Row Certification Lower Bound}
\label{app:cost-lowerbound}

Given $r \ge 1$ nonempty $\prec$-sorted runs with no cross-run order and a capacity-$m$ consistent oracle ($m \ge 2$), connect runs in a graph whenever a call ranks items from both. If at least two components remain at termination, shifting one component's ranks above all others preserves all observed permutations while changing $\pi_1$. Since one call touches at most $m$ runs, it merges at most $m - 1$ components, requiring $\lceil (r - 1) / (m - 1) \rceil$ calls—tightly matched by tournament elimination of up to $m$ run heads per call.
